\documentclass[11pt,a4paper]{article}
\usepackage[T1]{fontenc}
\usepackage{lmodern,amsmath,amssymb,amsthm,mathtools,booktabs,microtype}
\usepackage[margin=26mm]{geometry}
\usepackage[colorlinks=true,linkcolor=blue,citecolor=blue,urlcolor=blue]{hyperref}
\setlength{\emergencystretch}{2em}
\newtheorem{theorem}{Theorem}[section]
\newtheorem{lemma}[theorem]{Lemma}
\newtheorem{proposition}[theorem]{Proposition}
\newtheorem{corollary}[theorem]{Corollary}
\theoremstyle{remark}\newtheorem{remark}[theorem]{Remark}
\newcommand{\Q}{\mathbb Q}\newcommand{\Z}{\mathbb Z}\newcommand{\C}{\mathbb C}
\DeclareMathOperator{\Tr}{Tr}\DeclareMathOperator{\N}{N}\DeclareMathOperator{\Disc}{Disc}
\title{A single norm determinant for the 25- and 13-variable constructions}
\author{}\date{14 November 2026}
\begin{document}\maketitle
\begin{abstract}
We replace the product over all possible nonzero ranks in the previous
finite-algebra elimination by a single determinant. This is justified by
proving completeness with every conjugate denominator nonzero, not by
silently deleting the other strata. A free rank-four point algebra,
homogeneous elliptic differences, square extraction and a birational
rescaling of the pole parameters give further reductions. The resulting
degree bounds are $64674204453175296<2^{56}$ for 25 variables,
$155568385<2^{28}$ for 13 variables under weak $abc$ with exponent
$121/100$, and $126445825<2^{27}$ under exponent $16/15$.
The fully specified 27-variable construction has degree below $2^{59}$
and coefficient height below $2^{2^{72}}$, with the same certified
exponent $20483$. Counts include the input. The external arithmetic
premises remain explicit and are not certified by the finite checks.
No decrease in the number of variables is asserted.
\end{abstract}

\section{Statements, scope, and arithmetic inputs}
An integer polynomial $F(X,\boldsymbol Y)$ defines $\Z$ over $\Q$ when
its rational input projection is exactly $\Z$. All counts below include $X$.
Degree and support bounds are upper bounds, not expanded degrees or
monomial counts. Coefficient height is the largest absolute integer
coefficient after all fixed parameters have been substituted.

\begin{theorem}\label{thm:main}
Under the arithmetic premises specified below there are finite integer
polynomial specifications with the following bounds:
\begin{center}\begin{tabular}{lrrr}\toprule
Construction & Variables & Degree at most & Strict power bound\\\midrule
Inherited premises &25&64674204453175296&$2^{56}$\\
Fully specified, inherited premises&27&442502070853959680&$2^{59}$\\
Weak $abc$, exponent $121/100$&13&155568385&$2^{28}$\\
Weak $abc$, exponent $16/15$&13&126445825&$2^{27}$\\\bottomrule
\end{tabular}\end{center}
Each polynomial has integer content one. The 25-variable family defines
$\Z$ for every sufficiently large fixed odd exponent $L$; the sufficient
threshold is ineffective and need not equal the threshold for the older
squared-scale family. The 27-variable polynomial at $L=20483$ defines
$\Z$ and has
\[
 H(F)<2^{2^{72}}.
\]
In each 13-variable family every sufficiently large fixed odd shift
$B$ works, with an ineffective threshold. The displayed $L=3$ and $B=3$
computational fixtures are not certified defining members.
\end{theorem}

The inherited inputs are those stated precisely in the source papers
\cite{General,Explicit,Conditional}. They include the specified group
and index-height assertions for
\[
 K=\Q(\sqrt2),\quad D=(3-\sqrt2)^2=11-6\sqrt2,\quad
 E:y^2=x^3-Dx,
\]
\[
 P=\left(\frac{57-28\sqrt2}{8},\frac{-660+417\sqrt2}{32}\right),
 \qquad 2E(K)=\Z P;
\]
the prime-pattern theorem in its rational-contact form, with the
prescribed conjugate reduction orders, fixed finite divisibility and
avoidance, and real approximation in the prepared cyclic closure;
and the pole obstruction and its simultaneous elliptic realization.
The 25- and 27-variable routes also retain the full normalized modular
realization used for their contact-height bound. The trace-height
results and the explicit constants are exactly those used by their sources.
The two 13-variable routes instead assume the weak rational $abc$
inequality at the displayed exponent. Ordinary rational $abc$ implies
both; ordinary Szpiro supplies the first, as in the source.

In particular, the pole input used below is the \emph{two-curve
realization}, not just the assertion that one rational pole witness
exists. Its consequence, Zariski density in the two pole parameters,
is proved in \cite{Density} and recalled in Section~\ref{sec:density}.
This consequence uses an already retained arithmetic input. Assuming
only a black-box existence assertion would not justify our additional
nonvanishing requirements. None of these external arithmetic inputs is
claimed independently audited here. No new use of the Hexagon pole
replacement is needed for the present general-degree improvements.

\section{A free rank-four algebra for each point}\label{sec:points}
Retain $h_0,h_1\in\Q$, put $h=h_0+h_1\sqrt2$, and write
\[
 h^4-D=A+B\sqrt2,
\quad A=h_0^4+12h_0^2h_1^2+4h_1^4-11,
\quad B=4h_0^3h_1+8h_0h_1^3+6.
\]
The point source is $x=h^2$, $y=hj$, $j^2=A+B\sqrt2$.
Its rational solutions describe the nonzero multiples of $P$ by the
square-abscissa descent in the sources.

Over $\Z[A,B]$ use the free module with basis $1,U,V,W$ and the
following commutative multiplication table:
\begin{equation}\label{eq:table}
\begin{gathered}
 U^2=2A+2W,\qquad V^2=A-W,\qquad UV=2B,\\
 UW=AU-2BV,\qquad VW=BU-AV,\qquad W^2=A^2-2B^2.
\end{gathered}
\end{equation}
The letter $W$ in this table is local to the point block; detector
products below are denoted $\mathcal W$.

\begin{lemma}\label{lem:block}
The table defines an associative unital algebra. At every rational
$(h_0,h_1)$ it is etale of dimension four. Its rational characters
are exactly the pairs $j=u+v\sqrt2\in K$ with $j^2=A+B\sqrt2$,
under
\[
 U=2u,\qquad V=2v,\qquad W=u^2-2v^2.
\]
\end{lemma}
\begin{proof}
Substitution of $u^2+2v^2=A$, $2uv=B$ verifies all six relations.
After adjoining $\sqrt2$, the two values $u+\sqrt2v$ and
$u-\sqrt2v$ independently satisfy the square equations with radicands
$A+\sqrt2B$ and $A-\sqrt2B$. Where their product is nonzero these
four characters are distinct, and the four basis functions are
linearly independent. The table therefore agrees with multiplication
in an associative four-dimensional algebra over $\Q(A,B)$.
Associativity is a polynomial identity in $A,B$, so holds identically
over $\Z[A,B]$. The same character description proves etaleness when
$A^2-2B^2\ne0$.

At a rational pair, $h^4-D$ cannot vanish: its vanishing would imply
$\N_{K/\Q}(h)^4=\N(D)=49$, impossible for a rational norm. Since
$K$ is a field, $A^2-2B^2\ne0$. Finally the table implies
$u^2+2v^2=A$, $2uv=B$ and the displayed value of $W$, proving both
directions of the rational-character assertion.
\end{proof}

For root separation only, retain the former primitive element
\begin{equation}\label{eq:primitive}
 z=(2A+B)U-2(A+B)V.
\end{equation}
It satisfies
\[
 f(Z)=Z^4-4A'Z^2+8(B')^2,
\quad A'=6A^3-9AB^2-4B^3,
\quad B'=-4A^3+6AB^2+3B^3,
\]
with
\begin{equation}\label{eq:discriminant}
 \Disc(f)=2^{15}(B')^2(2A^2-B^2)^4(A^2-2B^2)^2.
\end{equation}
All factors are nonzero at a rational pair. For $B'$, if $B\ne0$
use Eisenstein at $3$ on $4T^3-6T-3$; if $B=0$ use $A\ne0$.
The other two factors are anisotropic quadratic forms over $\Q$.
The element $z$ generates the whole specialized point algebra:
with $m_0=2A+B$, $m_1=-A-B$, its inverse formulas are
\[
 U=\frac{m_0z-4m_1B'/z}{m_0^2-2m_1^2},\qquad
 V=\frac{-m_1z+2m_0B'/z}{m_0^2-2m_1^2}.
\]
The denominators are units, since $z\ne0$ and
$m_0^2-2m_1^2=2A^2-B^2\ne0$. Thus the four characters have
four distinct $z$ values. Importantly, these inverse formulas are
\emph{not} used to construct the arithmetic tests. Those use $U/2,V/2$
directly, avoiding the old high-degree denominators.

\subsection{Homogeneous differences and norms}
For points $(x,y)$ and $(x',y')$ on $E$, let
\begin{equation}\label{eq:difference}
\begin{split}
 d&=x-x',\\
 X_d&=(x+x')(xx'-D)+2yy',\\
 Y_d&=(yx'+y'x)d^2-(y+y')X_d.
\end{split}
\end{equation}
On $d\ne0$ their difference has coordinates
$x_d=X_d/d^2$, $y_d=Y_d/d^3$. Indeed the usual slope is
$(y+y')/d$, and substituting the two curve equations gives these
numerators. Hence
\[
 f(P-Q)=\frac{d^2}{X_d},\qquad
 z(P-Q)=-\frac{X_dd}{Y_d},\qquad f=1/x,\quad z=-x/y.
\]
For an original point, $z=-h/j$. Every field inversion retains its
norm denominator as a nonzero condition. We use
$\N(ab)=\N(a)\N(b)$, $\N(a^{-1})=\N(a)^{-1}$, and
$\N(j)=W$ from \eqref{eq:table}. These are identities in the finite
algebra, not numerical simplifications. The fourth-order formal
logarithm truncation is precisely $z$; no unused higher-order terms
or their denominators are introduced.

\section{The arithmetic sources and the smaller comparison scale}
We give the source formulas to fix the coefficient construction. Every
stated division is guarded, as are the contacts, pole inputs, and the
second pole roots. All scalar coordinates are rational.

For the 25-variable route retain nine point pairs, with point roles
\[
 A,S,T_1,Z,H,U_A,U_S,U_T,U_Z,
\]
one slope $\lambda\ne0$, and two pole pairs. Together with $X$ this
is 24 retained coordinates. Put $Q_1=A,Q_2=S,Q_3=T_1,Q_4=Z$,
$U_0=H$, $U_1=U_A,\ldots,U_4=U_Z$, and
\[
 V_0=H-P,\quad V_i=U_i-Q_i,\quad
 L_i=z(U_i)/z(H),\quad L'_i=\sigma(z(V_i)/z(V_0)).
\]
Here $\sigma$ acts on the displayed $K$ components. Define
\[
 \mathcal F(u,v,w)=v^2-u(u-2)vw+(u-2)^2w^2-1,
 \qquad \beta=216X^4+450459.
\]
The sixteen field tests are, in order,
\[
 (f(U_i))_{i=0}^4,\quad(\sigma f(V_i))_{i=0}^4,\quad
 (L_i-L'_i)_{i=1}^4,\quad L_1-\beta,\quad
 \mathcal F(L_2,L_3,L_4).
\]
Their orders are one except at positions $12,13,14,16$, which have
order three. The six scalar tests $S_k$ are the two rational components
of $f(H)$, $f(S-P)$, and $z(S-P)/z(H)$, with orders
$\kappa_k=(1,1,1,1,3,3)$.

For a fixed odd $L\ge3$, set
\begin{equation}\label{eq:rho}
\begin{gathered}
 \ell=76249,\quad
 r=\beta(\beta-1)(\beta-2401)(32\beta+343)(2\beta-1),\\
 C_L=17(\ell^{(L+1)/2}+1),\qquad
 \rho=\frac{102^2\ell^Lr}{(1+C_Lr)^2}.
\end{gathered}
\end{equation}
Let $T=\Tr\sum_i x(Q_i)$. The new scales are
\begin{equation}\label{eq:unsquared}
                         a_0=1,\qquad a_1=1+\rho T^2.
\end{equation}
This is the only change in the arithmetic source, as opposed to its
algebraic presentation. Put
\[
 \Theta=\{0,-1/2,-1,1/2399,-32/407\},\qquad
 \Theta_j=(4j+2+\Theta)\cup(4j+4+\Theta),\quad \gamma_j=-j-1,
\]
\[
 t_j=a_j\lambda,\quad \delta_j=t_j-\gamma_j,\quad
 F_j=\prod_{\theta\in\Theta_j}(t_j-\theta),\quad
 R_j=F_j/\delta_j^{10},\quad w_j=\rho R_j.
\]
With $\Phi_R(s)=Rs/(1+Rs)$ and $G_{j,k}=R_j^{\kappa_k}\Phi_{R_j}(S_k)$,
\begin{equation}\label{eq:detectors}
\begin{split}
 \mathcal W_j&=\prod_{i=1}^{16}
 \left(1+\frac{C_i^2}{\delta_j^2w_j^{2a_i-1}}\right)
 \prod_{k=1}^{6}\left(1+\frac{G_{j,k}^2}{w_j^{2\kappa_k-1}}\right),\\
 b_j&=w_j\frac{\N(\mathcal W_j)^2}{(\mathcal W_j)_0^2}
       =\frac{4w_j}{\Tr(\mathcal W_j^{-1})^2}.
\end{split}
\end{equation}
The fixed integer $a_i$ in the first product is the order of the $i$th
test, not the comparison scale.

For the explicit 27-variable route use three pole pairs and the same
nine point pairs. Write $\sum_i x(Q_i)=T_0+T_1'\sqrt2$, and retain
\[
 a_0=1,\quad a_1=(1+\rho T_0^2)^2,\quad
 a_2=(1+\rho(T_0+T_1')^2)^2.
\]
The contacts and detectors are the preceding ones for $j=0,1,2$.
These arithmetic scales are unchanged from \cite{Explicit}.

For the 13-variable routes retain the three point pairs $A,H,U$, one
slope, two pole pairs, and $X$, for 12 coordinates before elimination.
Put $\beta=X^2+B$, $T=\Tr x(A)$, and
\[
 a_0=1,\quad a_1=\beta^d+T^d,\qquad
 (\kappa,d)=(121/100,6)\ \hbox{or}\ (16/15,4).
\]
The six tests are
\[
 f(H),\ f(U),\ \sigma f(H-P),\ \sigma f(U-A),\
 z(U)/z(H)-\sigma(z(U-A)/z(H-P)),\ z(U)/z(H)-X.
\]
Use $\Theta=\{0,-1/2,-1/4,-2/5,-1,-1/3\}$, and for each scale put
\begin{equation}\label{eq:conditional}
 t=a\lambda,\quad \delta=t-1,\quad F=\prod_{\theta\in\Theta}(t-\theta),
 \quad \mathcal W=\prod_{i=1}^6(1+\delta^4C_i^2/F),\quad
 b=\frac{F\N(\mathcal W)^2}{4\delta^8\mathcal W_0^2}.
\end{equation}
The arithmetic source is unchanged from \cite{Conditional}.

\begin{proposition}\label{prop:unsquared}
Under the premises of \cite{General}, the source with scales
\eqref{eq:unsquared} has exactly the integer input projection for every
sufficiently large fixed odd $L$. Completeness holds for each integer
input and every fixed odd $L\ge3$.
\end{proposition}
\begin{proof}
We check the three uses of the old outer square. At an odd valuation of
$\rho$, $v(\rho T^2)$ is odd or is infinity. Thus $a_1$ is congruent
to one, or has negative valuation; there is no valuation-zero
cancellation in $1+\rho T^2$. The disjoint contact-and-pole sets give
a unit among $R_0,R_1$, exactly as in the scalar-odd lemma of the source.
The proof of that lemma uses only this alternative, not parity of
$v(a_1)$. All local test-extraction and residual divisibilities remain
unchanged.

For the height step the new inequality is
\[
 h(a_1)\ge2h(T)-h(\rho)-\log2,
 \qquad h(a_1)\le h(t_0)+h(t_1).
\]
Only the fixed constants change. In the notation of the source, with
$c=|a-1|$ and $B_0$ the maximum absolute root index, the same bounds give
\[
 h(\beta)\le H_1c^{7/9},\quad c\ge\ell^{2L+1},\quad
 B_0^2\le C|Nc|^\eta+Cc^{7/9}+C\log c+J,
\]
where $\eta=7/18$ for a nonzero Lucas residual
$N=\mathcal F(n_2,n_3,n_4)$ and $\eta=7/6$ for a nonzero input
residual $N=u-n_1v$, $\beta=u/v$ reduced. Constants are independent
of the input, witnesses, and $L$. If the Lucas residual is nonzero,
$|N|\le5B_0^4$ and $c\le2B_0$ give an upper exponent $35/18<2$;
hence $B_0$ is bounded, contradicting the lower bound on $c$ for large $L$.
Its vanishing gives the source's integer Lucas growth
$\log B_0>c-4$. If the input residual is nonzero, use
$|N|\le2e^{h(\beta)}B_0$ and absorb the lower-order terms. Taking
logarithms yields
\[
 \log B_0\le C'+\frac75 h(\beta)+\frac75\log c=o(c),
\]
a contradiction. Thus $\beta=n_1\in\Z$. Integrality of
$216X^4+450459$ forces integrality of $X$, since a reduced denominator
$v$ would satisfy $v^4\mid216$.

For completeness, after the root points have been fixed the scales are
positive rational numbers. The prime-pattern application permits these
fixed rational coefficients. For $\lambda=u/v$ its squareclass identity is
\[
 R_j=\left(\frac{a_j}{v\delta_j}\right)^{10}
       \prod_{\theta\in\Theta_j}\bigl(u-(\theta/a_j)v\bigr).
\]
The exponent ten already makes the prefactor a square: $a_j$ itself
need not be a square. The prescribed primes and real and finite-place
limits $R_j\to1$, $\delta_j^{-1}\to0$ are unchanged. The extra
common-index divisibility at the comparison primes, the scalar
normalization argument, and the preparation of both $b_j$ therefore
apply verbatim. In particular the prepared range is still $(0,17)$.
This checks every use of the deleted outer square.
\end{proof}

\section{Rescaling the two retained pole parameters}\label{sec:poles}
Write a nonzero pole input as $b=d t^2$, with rational nonzero $d,t$.
We make this decomposition by flattening products and powers in the
fixed numerator and denominator circuits, putting even exponents into
$t^2$. This is syntactic square extraction, not an irreducible
factorization or a squarefree assertion about polynomial factors.
Use retained parameters $\tau,s$ and set
\begin{equation}\label{eq:newpole}
 e=\frac{\tau^2-17d}{\tau^2+51d},\qquad
 q^2=17d\bigl(s^2-17e(17-b)\bigr)
                  \bigl(s^2+51e(17-b)\bigr),\qquad q\ne0.
\end{equation}
The denominator is guarded nonzero.

\begin{lemma}\label{lem:rescale}
On $bt\ne0$, \eqref{eq:newpole} is equivalent, by invertible rational
coordinate changes, to the former one-conic-chart pole system. Its
retained parameters and second ordinate are
\[
 \tau_{\rm old}=t\tau,\qquad s_{\rm old}=bs,
 \qquad q_{\rm old}=\frac{bt}{17}q.
\]
\end{lemma}
\begin{proof}
Put $B_0=17b$. The old chart gives
$e=(\tau_{\rm old}^2-B_0)/(\tau_{\rm old}^2+3B_0)$,
which becomes the displayed $e$. Since
$B_0e(B_0-b^2)=17b^2e(17-b)$, the old radicand becomes
\[
 \frac{b^3}{17}
       (s^2-17e(17-b))(s^2+51e(17-b)).
\]
Dividing by $(bt/17)^2=b^3/(289d)$ gives \eqref{eq:newpole}.
All transformations are invertible and preserve nonzero ordinates.
The first ordinate reconstructs as
$-68dt\tau/(\tau^2+51d)$. Thus the pole obstruction is unchanged.
\end{proof}

If the last radicand is $N/D_0$, clear it by a finite square root
$Z^2=ND_0$, $Z\ne0$, with $D_0\ne0$. Flatten multiplication and powers
again to write $ND_0=S_0^2c$. Replace $Z$ by $S_0q_0$, and use
$q_0^2=c$. All factors of $S_0$, of $c$, and of the source nonzero
conditions will be included in the support product defined below.
No unrestricted coordinate is added.

\section{Completeness with all conjugate denominators nonzero}\label{sec:density}
This section is essential: the new polynomial deliberately excludes
some retained base tuples which the former polynomial allowed.

\begin{lemma}[Density of the prepared pole parameters]\label{lem:density}
Under the retained simultaneous elliptic realization, a prepared
nonzero input $b$ has a Zariski-dense set of rational witness pairs in
the old parameter plane $(\tau_{\rm old},s_{\rm old})$. The same is true
in the rescaled plane $(\tau,s)$.
\end{lemma}
\begin{proof}
We recall the argument of \cite{Density}. Put $B_0=17b$ and
$V=B_0(B_0-b^2)=\theta c^2>0$. The realization supplies positive
$n$ and infinite-order points on
\[
 E_n:y^2=x(x-n)(x+3n),\qquad E_{n\theta},
\]
with abscissas in $B_0\Q^{\times2}$. The non-torsion assertion follows
from the source's fresh-prime condition and the torsion lemma in
\cite{Density}: for nonsquare $a$, $E_a(\Q)$ has only two-torsion.
That lemma follows from Mazur's theorem, absence of roots modulo five
of $3z^4+8z^3-18z^2-9$, and the doubling test at the three two-torsion
points.

Independently vary positive odd multiples and write their abscissas
as $B_0u^2,B_0v^2$. The two-isogeny squareclass homomorphism preserves
these classes. The resulting witnesses have
\[
 e=n/(B_0u^2),\quad p=y_1/(B_0u^3),\quad
 s_{\rm old}=cv/u,\quad q_{\rm old}=c^2y_2/(B_0^2u^2v).
\]
Here $p,q_{\rm old},s_{\rm old}$ are nonzero and $0<e<1$.
Varying the first multiple gives infinitely many $e$, hence infinitely
many $\tau_{\rm old}=p/(e-1)$. For each fixed first multiple, varying
the second gives infinitely many $s_{\rm old}$. A polynomial vanishing
on all these pairs vanishes first as a polynomial in $s_{\rm old}$,
then coefficientwise at infinitely many $\tau_{\rm old}$; it is zero.
This argument works over $\overline\Q$ as well. The rescaling of
Lemma~\ref{lem:rescale} is an invertible diagonal linear map for fixed
$b,t$, and preserves density.
\end{proof}

\begin{lemma}\label{lem:avoidance}
Let $B_0$ be a finite etale $\Q$-algebra. Suppose the elements $b,d,t$
are units and $b=dt^2$. In the rescaled pole formulas,
$\N_{B_0}(\tau^2+51d)$ is a nonzero monic polynomial of degree
$2\dim B_0$. After excluding it, the norm of the radicand in
\eqref{eq:newpole}, as a polynomial in $s$, has degree
$4\dim B_0$ and nonzero leading coefficient $\N_{B_0}(17d)$.
Consequently, if the actual rational character sends $b$ to a prepared
input, a rational pole witness can be chosen so that every native
pole denominator and every conjugate radicand is nonzero.
\end{lemma}
\begin{proof}
The first norm is the determinant of $\tau^2I+51M_d$. The second
multiplication matrix has leading coefficient $17M_d$ in degree four.
Both determinants are therefore nonzero polynomials. Clear the finitely
many already nonzero scalar denominators. Every numerator or denominator
used in forming the pole is a product of fixed unit factors, the conic
denominator and the two monic quadratic factors in $s$; these give
nonzero specialized polynomials at every geometric character. Their
finitely many norm conditions can be avoided simultaneously by
Lemma~\ref{lem:density}. Nonzero products imply nonvanishing of every
syntactically extracted factor. Independent prepared poles are handled
successively, or by the product of the dense witness sets.
\end{proof}

\begin{proposition}\label{prop:fullunits}
For every integer input in the completeness range of each arithmetic
source, there is a rational source witness for which every nonzero
condition and every pole radicand is nonzero at \emph{every} geometric
character of the point algebra. For the 13-variable sources this is
asserted for nonzero integers; zero is included by the final factor $X$.
\end{proposition}
\begin{proof}
Follow the order of choices in the source completeness proofs. For the
general routes choose the scalar label $a=c+1>1$ and the four nonzero
root labels first. For the conditional routes fix the root label $n=X\ne0$.
After these choices the scales and all root abscissas are fixed rational
quantities. Temporarily take $H=[h]P$ and $U_i=[hn_i]P$.

At an actual rational point pair, $j$ has two real embeddings, both
nonzero. Each geometric character of the point block amounts to choosing
their two signs independently. It therefore changes, independently at
the real places, each ordinate to one of its two signs, while leaving
its abscissa fixed. All point-algebra characters are real at these
completeness tuples. A difference on any such branch is one of
$[\pm h\pm1]P$ or $[\pm hn_i\pm n_i]P$ at the respective embedding.
The fixed scalar difference is $[a-1]P$ or $[-a-1]P$.
These points are neither the identity nor two-torsion when $h>2$,
$n_i\ne0$, and $a>1$. Equal-abscissa denominators are also nonzero.
Thus the point, difference, reciprocal and ratio denominators on all
sign branches can be kept nonzero in a sufficiently small real
neighborhood of a prepared return. There are finitely many branches. No continuously varying rational
square root of an abscissa is required here: the guard functions in
question are functions of the points and their sign choices; the
additional $h,j$ norm factors are already nonzero.
The elementary $h,j$ denominators in the new formulas are nonzero
there as well; their sign choices only permute these branches.

The conditional source has no scalar normalizers. For a general source
consider the additional denominators $1+R_j(\lambda)D(H)$, where $D$
is any scalar component on any sign branch. On a connected real
coordinate neighborhood in the prepared cyclic closure these are real
analytic. If $D$ is a constant $d\ne0$, the equation
$R_j(\lambda)=-1/d$ has only finitely many solutions, since the contact
function $R_j$ is nonconstant. Increase the lower bound on the positive
slope beyond all such solutions before prime selection. Constants zero
impose no condition. If $D$ is nonconstant, for each subsequently fixed
allowed slope the zero set of $1+R_j(\lambda)D(H)$ has empty interior.
Deleting the finitely many such zero sets leaves a nonempty open subset
of the prepared neighborhood.

This deletion is compatible with the arithmetic order of choices.
The initial congruence modulus annihilates the finite component group
of the compact cyclic closure. Its identity component is a real torus;
multiplication by any positive integer is surjective. Hence each
compatible refinement of the common-index progression, including the
conjugate reduction-order conditions and the additional selected-prime
divisibility of the general source, is dense in the same component.
After prime selection it therefore meets the remaining open subset,
with $h>2$. All previously prepared real inequalities and finite-place
conditions persist. Only avoidance of a finite number of exceptional
slope values was added before the prime-pattern application.

For sufficiently large positive slope all $R_j,w_j$ in the general
source are positive, and all weights $\delta^4/F$ in the conditional
source are positive. Each detector factor at either real embedding,
on every point sign branch, is one plus a nonnegative square times a
positive coefficient. Thus $\mathcal W_+,\mathcal W_->0$,
$\mathcal W_0=(\mathcal W_++\mathcal W_-)/2>0$, and
$\N(\mathcal W)>0$ on every character. Both forms of $b_j$ are therefore
nonzero on every character: they are units of the point algebra.
Their syntactic square decompositions have unit factors, since both
the numerator and the denominator are already nonzero there.

The actual character still has the prepared pole inputs supplied by the
original completeness argument (or Proposition~\ref{prop:unsquared}).
Apply Lemma~\ref{lem:avoidance} independently to these inputs. The
remaining pole parameters then make all conic denominators, all cleared
radicands and their extracted factors units. These are all denominator
types in the displayed rational source: field norms, scalar normalization
denominators, fixed coefficients, contacts, and the pole denominators.
No residual equality or further denominator is omitted. This proves the
proposition.
\end{proof}

\section{The single determinant}\label{sec:determinant}
Let $R=\Z[x_1,\ldots,x_a]$ be the ring of retained coordinates. Let
$B_0$ be the tensor product of the $p$ point blocks; it is free of rank
$b=4^p$. After the square extractions let
\[
 A_0=B_0[q_1,\ldots,q_s]/(q_j^2-c_j),\qquad n=2^sb.
\]
Let $T_*$ be the product, without repetitions, of all nonconstant
syntactic factors in the source nonzero conditions and in all cleared
radicands, after the pole generators themselves have been removed.
Nonzero integer constants are omitted; a fixed parameter is omitted
only if it is a known nonzero constant. The supplied specifications
retain all parameter atoms. Every factor lies in $B_0$.
Define the scalar polynomial
\begin{equation}\label{eq:normguard}
                         \Delta_* = \det_{B_0}M_{T_*}.
\end{equation}
If it is nonzero at a rational base tuple, every source denominator,
every square multiplier and every $c_j$ is a unit. If it vanishes, the
final polynomial below is identically one in the extra coordinate.

For each point use $z_i,f_i$ from \eqref{eq:primitive}. Write
$f_i(Z)=Z^4+\sum_{k<4}f_{ik}Z^k$, and define, solely by integer
polynomial and matrix operations,
\begin{equation}\label{eq:separator}
\begin{gathered}
 Z_0=2+\sum_{i,k<4}f_{ik}^2,\qquad
 S=b\left(1+\sum_{j=1}^s\sum_{u,v=1}^b(M_{c_j}^{B_0})_{uv}^2\right),\\
 \delta_i=\Disc(f_i),\quad D_0=\prod_i\delta_i^2,\quad
 K_0=(2Z_0)^5\prod_i(1+\delta_i^2),\quad C_0=D_0+4Z_0K_0,\\
 J=4S,\quad Q=2SJ^{s-1},\quad
 \Xi=\sum_{j=1}^s J^{s-j}\left(\prod_{k<j}c_k\right)q_j,\\
 \Omega=Q\left(C_0^{p+1}+\sum_{i=1}^pD_0^iC_0^{p+1-i}z_i\right)
                      +D_0^{p+1}\Xi.
\end{gathered}
\end{equation}
The matrices use the tensor basis $1,U_i,V_i,W_i$ for each point,
followed by the square exponents $0,1$ in the pole generators.

\begin{lemma}\label{lem:separator}
At a rational base tuple with $\Delta_*\ne0$, the values of $\Omega$
on the $n$ geometric characters of $A_0$ are distinct and nonzero.
\end{lemma}
\begin{proof}
We include the separation argument from \cite{Triangular} to specify
its use in the new basis. Cauchy's bound gives $|z_i|\le Z_0$.
The spectral-radius bound gives $|c_j|\le S$. Put
$\epsilon=D_0/K_0$ and $L=C_0/D_0=1+4Z_0/\epsilon$.
All discriminants are nonzero, and $0<\epsilon\le1$.
If $d$ is one root difference of $f_i$, its discriminant gives
$|\delta_i|\le |d|^2(2Z_0)^{10}$. Since
$t^2/(1+t^2)\le\sqrt t$ for $t>0$, one obtains $|d|\ge\epsilon$.

For nonzero $c_j$ the values
\[
 \phi=\frac1{2S}\sum_{j=1}^s(4S)^{1-j}
                       \left(\prod_{k<j}c_k\right)q_j
\]
are distinct on the $2^s$ square-root choices and have absolute value
at most $2/3$. For two choices first differing at $j$, divide out their
common earlier product. The leading difference is $2q_j$; the ratio
of the remaining absolute sum to its magnitude is at most
$(|q_j|/\sqrt S)\sum_{r\ge1}4^{-r}\le1/3$. This uses no lower
bound for the nonzero radicands and permits nonreal values.

Divide $\Omega$ by the nonzero rational scalar $QD_0^{p+1}$. The result is
\[
 L^{p+1}+\sum_{i=1}^pL^{p+1-i}z_i+\phi.
\]
At the first differing point coordinate, the leading difference has
magnitude at least $\epsilon L^k$, while the tail is strictly less
than $2Z_0L^k/(L-1)=\epsilon L^k/2$. If no point coordinate differs,
the preceding quadratic argument applies. Finally the distance from
$L^{p+1}$ is less than
$Z_0L^{p+1}/(L-1)\le L^{p+1}/4$. Thus the values are nonzero.
\end{proof}

Define
\begin{equation}\label{eq:finaldet}
                       P(\boldsymbol x,U)=
                       \det_{A_0}(I-U\Delta_*M_\Omega).
\end{equation}
For the conditional sources set $F=XP$; otherwise set $F=P$.

\begin{theorem}\label{thm:norm}
If $\Delta_*=0$ at a rational base tuple, $P(U)=1$ for every rational
$U$. If $\Delta_*\ne0$, $P$ has a rational zero if and only if the
corresponding guarded rational source has a rational character.
Together with Proposition~\ref{prop:fullunits}, this proves that the
new polynomials have the asserted integer input projections.
\end{theorem}
\begin{proof}
The first assertion follows directly from the identity matrix in
\eqref{eq:finaldet}; it is true even if the square extensions have
nilpotents. For the second, $A_0$ is etale and its character values of
$\Delta_*\Omega$ are distinct nonzero numbers by Lemma~\ref{lem:separator}.
A rational zero $u$ is necessarily nonzero. The kernel over $\Q$ of
$I-u\Delta_*M_\Omega$ is one-dimensional. Every multiplication
operator preserves this kernel, so acts on it by a rational scalar.
These scalars define a unital $\Q$-algebra homomorphism $A_0\to\Q$.
All source guards and square multipliers are units, so the rational
reconstruction gives the guarded source, including its nonzero pole
roots. Conversely, a rational character supplies
$u=1/(\Delta_*\Omega)$ and hence a rational zero.

Thus every nonzero-input rational zero belongs to the arithmetic source;
its soundness applies. Completeness for each required integer is supplied
by Proposition~\ref{prop:fullunits}. The conditional zero input is
included by $XP$. Finally $P(\boldsymbol x,0)=1$ identically, proving
integer content one for $P$ and also for $XP$.
\end{proof}

\begin{remark}
We do not claim equality of the old and new rational projections to
\emph{all} retained coordinates. The new construction excludes any
base tuple with a zero conjugate denominator, even when another
character is valid. The additional completeness proof is what preserves
the projection to the distinguished input. No general inference from
Zariski density to the existence of a point on a prescribed hypersurface
is used: only avoidance of finitely many nonzero polynomials is required.
\end{remark}

\section{Degree and coefficient bounds}\label{sec:bounds}
Base coordinates have weight one. In a point block use weights
$0,2,2,4$ for $1,U,V,W$. The multiplication table does not increase
weight. For coefficient $\ell^1$ norms use diagonal basis scales
$1,16,16,64$. The exact point coefficients have
\[
 \|A\|_1=28,\qquad \|B\|_1=18.
\]
The weighted column norms of multiplication by $U,V,W$ are at most
$16,16,64$, respectively. For example the last column for $U$ is
bounded by $(28+2\cdot18)16/64=16$; the second is
$56/16+2\cdot64/16<16$. The other columns follow immediately from
\eqref{eq:table}. The tensor basis therefore has maximum degree weight
$4p$ and maximum logarithmic scale $6p$.

All circuit bounds use the following rules. Degree weights add under
multiplication and take their maximum under addition. Upper bounds for
$\log_2$ of operator coefficient norms add under multiplication; an
addition of $k$ terms costs at most $\lceil\log_2 k\rceil$ beyond
the maximum. A fixed integer has weight zero. A pole radicand with
bounds $(d_j,\ell_j)$ gives the root weight
$v_j=\lceil d_j/2\rceil$ and root logarithmic norm bound
$g_j=\lceil\ell_j/2\rceil+1$. These rules follow from the weighted
companion matrix and do not require expansion of the production matrices.

Let $(t,\tau)$ bound the weight and logarithmic operator norm of $T_*$.
Then
\[
 \deg\Delta_*\le bt,\qquad \log_2\|\Delta_*\|_1\le b\tau.
\]
For the norm bound, expand the determinant by columns: its coefficient
norm is at most the product of column sums. Diagonal basis scales cancel
in each determinant term. This also proves the degree bound by cancellation
of row and column weights.

For completeness we give the exact remaining recurrences. Write
$\ell(k)=\lceil\log_2k\rceil$. Degree and logarithmic norm bounds for
$Z_0,D_0,C_0,S,Q$ are
\begin{align*}
 z_d&=48,& z_l&=78+\ell(2+2p),\\
 D_d&=144p,& D_l&=234p,\\
 C_d&=288+D_d,&C_l&=5(1+z_l)+p+D_l+z_l+3,\\
 S_d&=2\max_j(d_j+4p),&S_l&=\ell(b)+\ell(1+sb^2)+2\max_j(\ell_j+6p),\\
 Q_d&=sS_d,&Q_l&=sS_l+2s-1.
\end{align*}
Indeed $\|A'\|_1\le236688<2^{18}$ and
$\|B'\|_1\le159736<2^{18}$; the two nonzero coefficients of $f_i$
have logarithmic bounds $20,39$. Also
$\|2A^2-B^2\|_1\le1892<2^{11}$ and
$\|A^2-2B^2\|_1\le1432<2^{11}$, so \eqref{eq:discriminant}
gives a discriminant logarithmic bound $117$. The primitive element
$z_i$ has weight six and norm at most $2656<2^{12}$.

Put
\begin{align*}
 P_d&=Q_d+\max\bigl((p+1)C_d,
                  \max_i(iD_d+(p+1-i)C_d+6)\bigr),\\
 P_l&=Q_l+\ell(p+1)+\max\bigl((p+1)C_l,
                  \max_i(iD_l+(p+1-i)C_l+12)\bigr),\\
 X_d&=\max_j\bigl((s-j)S_d+\sum_{k<j}d_k+v_j\bigr),\\
 X_l&=\ell(s)+\max_j\bigl((s-j)(S_l+2)+\sum_{k<j}\ell_k+g_j\bigr),\\
 a_d&=\max(P_d,(p+1)D_d+X_d),\\
 a_l&=1+\max(P_l,(p+1)D_l+X_l).
\end{align*}
They bound $\Omega$. Finally
\begin{equation}\label{eq:finalbounds}
 \deg P\le n(bt+a_d+1),\qquad
 \log_2\|P\|_1\le n(b\tau+a_l+1),\qquad \deg_U P\le n.
\end{equation}
The conditional input factor adds one to degree and does not change
coefficient norm or support size. If $V$ is the total variable count and
$D_P=n(bt+a_d+1)$, a support upper bound is
\[
 \binom{D_P+V}{V}-\binom{D_P-n+V-1}{V}.
\]
This counts the monomials of total degree at most $D_P$ whose $U$ exponent
is at most $n$; it remains valid after multiplying all monomials by $X$.

\begin{center}\begin{tabular}{lrrrr}\toprule
Construction & $p$ & $s$ & $t$ & $a_d$\\\midrule
25 variables &9&2&235281&625256\\
27 variables &9&3&804900&1724364\\
13, exponent $121/100$&3&2&9127&23560\\
13, exponent $16/15$&3&2&7427&18600\\\bottomrule
\end{tabular}\end{center}
Substitution in \eqref{eq:finalbounds} gives Theorem~\ref{thm:main}.
The support upper bounds are respectively $2^{1282},2^{1457},2^{306}$
and $2^{303}$, all strict. At the certified value $L=20483$, the
27-variable logarithmic coefficient bound is
\[
 \log_2\|F\|_1\le3707461411940344004608<2^{72}.
\]
The 25-variable and conditional coefficient certificates refer only to
their stored small-parameter fixtures. Their degree bounds are uniform
in the permitted fixed parameters: integer coefficient leaves have
weight zero, and the multiplicative exponents of nonconstant factors
in this fixed construction do not change with those parameters.

\section{Reproducible specification and remaining questions}
The complete native source generator is
\nolinkurl{computation/few_variables/unit_norm_v28/build_sources.py}.
The four specifications under its \texttt{specifications} directory
fix the generator arguments, the canonical source hashes, the point
blocks, the exact syntactic square extractions, every support factor,
and the finite determinant recipe. The complete integer circuits can
be exported by the builder's \texttt{--native-sources} option. No
unrecorded choice of a favorable root or a favorable rank is part of
the polynomial definition.

The independent verification separates polynomial identities and
exact factor certificates from finite-field tests and small examples.
The latter test the implementation and are not proofs of any infinite
arithmetic input. Deliberate corruptions test whether the checks reject
incorrect multiplication, square extraction, bounds, and boundary handling.
The checks include three complete eight-dimensional examples of the
separator, a partially vanishing guard, fourteen fixed-parameter
reconstructions, and 180 nonsingular finite-field formula evaluations.
There are eight altered-degree or altered-sign rejection tests, together
with separate controls for the multiplication table, pole rescaling,
nonseparation and omission of the scalar guard norm.
The earlier source manuscripts, code and byte-pinned data remain unchanged.

The variable counts 25 and 13 have not been reduced. The point and pole
parameters are still needed by the completeness proof. In particular,
density in a two-dimensional pole-parameter plane does not justify
identifying two independent parameters, fixing one arbitrarily, or
sharing parameters belonging to different pole inputs. The next reduction
must establish a suitable arithmetic realization for such a restriction,
or replace a retained elliptic point role, rather than infer a dimension
reduction from the present avoidance argument.

\begin{thebibliography}{9}
\bibitem{General}\emph{A defining polynomial in twenty-five variables:
combined scalar and comparison conditions}, project source, 10 November 2026,
\nolinkurl{papers/few_variables25/few_variables25_merged_20261110_v19.tex}.
\bibitem{Explicit}\emph{An Explicit Defining Polynomial in Twenty-Seven Variables}, project source, 10 November 2026,
\nolinkurl{papers/few_variables27/explicit27_20261110_v19.tex}.
\bibitem{Conditional}\emph{Thirteen Variables from Direct Elliptic Index
Comparison}, project source, 10 November 2026,
\nolinkurl{papers/conditional_few_variables/direct_input13_20261110_v22.tex}.
\bibitem{Density}\emph{Zariski density of the prepared pole witnesses},
project source, 9 November 2026,
\nolinkurl{papers/pole_parameters/pole_witness_density_20261109_v14.tex}.
\bibitem{Triangular}\emph{Separating Quadratic Roots in Finite Algebras and Smaller Degrees for Diophantine Presentations},
project companion, 13 November 2026,
\nolinkurl{papers/finite_algebra/triangular_separation_20261113_v24.tex}.
\end{thebibliography}
\end{document}
